\newcommand{\figuraHidrometroProblema}{%
\begin{figure}[H]
    \centering
    \begin{tikzpicture}[>=Latex,font=\small,line cap=round,line join=round]
        \draw[very thick] (2.0,5.6)--(2.0,0.4)--(9.8,0.4)--(9.8,5.6);
        \fill[blue!12] (2.08,0.48) rectangle (9.72,3.82);
        \draw[very thick,blue!65!black] (2.0,3.82)--(9.8,3.82);
        \node[blue!65!black] at (3.0,2.75) {agua};

        \fill[gray!25] (5.9,1.55) ellipse (0.78 and 1.05);
        \draw[very thick] (5.9,1.55) ellipse (0.78 and 1.05);
        \draw[very thick] (5.68,2.40)--(5.68,5.25)--(6.12,5.25)--(6.12,2.40);
        \begin{scope}
            \clip (5.9,1.55) ellipse (0.78 and 1.05);
            \fill[gray!65] (5.1,0.50) rectangle (6.7,1.30);
        \end{scope}
        \draw[very thick] (5.9,1.55) ellipse (0.78 and 1.05);
        \node[white,font=\scriptsize] at (5.9,1.12) {Hg};
        \foreach \y in {2.70,3.05,3.40,3.75,4.10,4.45,4.80}
            \draw (6.12,\y)--(6.35,\y);

        \draw[<->,thick] (6.95,2.40)--(6.95,5.25)
            node[midway,right] {$12\ \mathrm{cm}$};
        \draw[very thick,red!75!black] (5.52,3.82)--(6.28,3.82);
        \fill[red!75!black] (5.90,3.82) circle (0.055);
        \node[red!75!black,anchor=west] at (7.70,3.48)
            {punto medio};
        \node[align=left] at (3.75,4.62)
            {$m_{\mathrm{vidrio}}=0.01\ \mathrm{kg}$\\
             masa de Hg desconocida};
    \end{tikzpicture}
    \caption{Hidrómetro del problema. El lastre de mercurio se coloca en el
    bulbo inferior y debe elegirse para que, en agua, la superficie corte el
    punto medio del vástago de $12\,\mathrm{cm}$.}
    \label{fig:hidrometro-problema}
\end{figure}%
}